\documentclass[compacto,10pt]{pubciencias-notas}

% -- Paquetes adicionales
\usepackage{comandosEPN}

% -- Datos del libro
\autor{Departamento de Formación Básica}
\materia{Álgebra Lineal}
\nota{Semana 11 Clase 01}
\fecha{Semestre 2025-A}

\logouno[1.1cm]{Logo_EPN.eps}
\logodos[1.7cm]{Logo_DFB.eps}

\definecolor{colordef}{cmyk}{0.81,0.62,0.00,0.22}

\begin{document}

\encabezado
\vspace*{-8mm}

\setcounter{section}{4}
\section{Espacios con producto interno}
\subsection{Producto interno}

\begin{defi}[Producto interno]
    Sea $(E,+,\cdot,\K)$ un espacio vectorial. Un producto interno sobre $E$ es una función
    \[
        \funcion{\langle  \cdot, \cdot \rangle }{E\times E}{\K}{(u,v)}{\langle  u, v \rangle }
    \]
    tal que cumple:
    \begin{enumerate}
        \item 
            $\langle  v, v \rangle  \geq 0$ para todo $ v \in E$;
        \item 
            $\langle  v,v \rangle  = 0$ si y solo si $ v = 0$;
        \item 
            $\langle  u + v, w \rangle  = \langle  u,w \rangle  + \langle  v,w\rangle  $ para todo $u,v,w \in E$;
        \item
            $\langle  \alpha v,w \rangle  = \alpha \langle  v,w \rangle $ para todo $v,w \in E$ y $\alpha\in\K$.
        \item
            $\langle  v,w \rangle  = \overline{\langle w, v \rangle}$ para todo $v,w \in E$.
    \end{enumerate}
\end{defi}


\begin{advertencia}
    En caso de tratarse de un espacio vectorial sobre el campo de los números reales, la propiedad 5 equivale a
    \begin{enumerate}[start=5]
        \item
            $\langle  v,w \rangle  = \langle  w, v \rangle $ para todo $v,w \in E$.
    \end{enumerate}
\end{advertencia}


\begin{advertencia}
    Se puede demostrar utilizando la propiedad 4 que
    \begin{itemize}
        \item
            si $v = 0$, entonces $\langle  v,v \rangle  = 0$;
    \end{itemize}
    por lo cual, podríamos cambiar la propiedad 2 por
    \begin{enumerate}[start=2]
        \item
            si $\langle  v,v \rangle  = 0$, entonces $ v = 0$
    \end{enumerate}
\end{advertencia}

\begin{advertencia}
    Si se define un producto interno sobre un espacio vectorial $E$, a este se lo denomina espacio con producto interno o prehilbertiano.
\end{advertencia}

%\begin{ejem}\hspace{0pt} En el espacio vectorial $\R^3$ definimos a la función 
%\[
%    \funcion{\langle \cdot,\cdot\rangle }{\R^{3} \times \R^{3}}{\R}{(u, v)}{u_1 v_1}
%\]
%donde $u = (u_1, u_2, u_3)$ y $v  = (v_1,v_2,v_3)$. Tomando $u = (0,2,3)$, se tiene que
%\[
%    \langle u,u\rangle  = \langle (0,2,3),(0,2,3)\rangle  = 0,
%\]
%sin embargo, $u \neq (0,0,0)$, por tanto, $\langle \cdot,\cdot\rangle $ no es un producto interno.
%\end{ejem}

%\begin{ejem}\hspace{0pt} En el espacio vectorial $\R_1[x]$ definimos la función $$\func{\langle \cdot,\cdot\rangle }{\R_1[x] \times \R_1[x]}{\R}$$ por
%\[
 %   \langle  p, q \rangle  = a_0 b_0 + 2 a_1b_1 
%\]
%para todo $p(x) = a_0 + a_1 x\in\R_1[x]$ y todo $q(x) = b_0 + b_1x\in\R_1[x]$. La función $\langle \cdot,\cdot\rangle $ es un producto interno.

%\begin{proof}[Demostración.]\hspace{0pt}
%Sean $p,q,r\in\R_{1}[x]$ y sea $\alpha\in\R$.
%\begin{enumerate}
%\item Tenemos que $p(x) = a_{0} + a_{1}x$, y que
%\[
 %   \langle p,p\rangle  = a^{2}_{0} + 2a^{2}_{1} \geq 0.
%\]
%\item Se tiene que
%\begin{align*}
 %   \langle p,p\rangle  = 0
  %  &\Longleftrightarrow 
   % a^{2}_{0} + 2a^{2}_{1} = 0.
%\end{align*}
%Lo cual se verifica si y solo si $a_{0} = a_{1} = 0$, o equivalentemente $p(x) = 0 + 0x$.
%\item Notemos que $q(x) = b_{0} + b_{1}x$ y $r(x) = c_{0} + c_{1}x$, entonces
%\begin{align*}
 %   \langle p + q, r\rangle  
  %  &= (a_{0} + b_{0})c_{0} + 2(a_{1} + b_{1})c_{1}\\
   % &= (a_{0}c_{0} + 2a_{1}c_{1}) + (b_{0}c_{0} + 2b_{1}c_{1})\\
    %&= \langle p, r\rangle   + \langle q, r\rangle .
%\end{align*}
%\item Tenemos que
%\begin{align*}
 %   \langle \alpha p, q\rangle  = (\alpha a_{0})b_{0} + 2(\alpha a_{1})b_{1} = \alpha(a_{0}b_{0} + 2a_{1}b_{1}) = \alpha\langle p,q\rangle .
%\end{align*}
%\item Se tiene que
%\begin{align*}
 %   \langle p,q\rangle  = a_{0}b_{0} + 2a_{1}b_{1} = b_{0}a_{0} + 2b_{1}a_{1} = \langle q,p\rangle .
%\end{align*}
%\end{enumerate}
%Por 1, 2, 3, 4 y 5 concluimos que $\langle \cdot,\cdot\rangle $ es un producto interno.\qedhere
%\end{proof}
%\end{ejem}

%\begin{ejer}\hspace{0pt} En el espacio vectorial $\Mat{2}{2}$ se define a la función
%\[
 %   \func{\langle \cdot,\cdot\rangle }{\Mat{2}{2}\times\Mat{2}{2}}{\R}
%\]
%por 
%\[
 %   \langle A,B\rangle  = a_{11}b_{11} + 2a_{12}b_{12} + 2a_{21}b_{21} + a_{22}b_{22}
%\]
%para toda $A = \begin{pmatrix}a_{11}&a_{12}\\a_{21}&a_{22}\end{pmatrix}\in\Mat{2}{2}$ y toda $B = \begin{pmatrix}b_{11}&b_{12}\\b_{21}&b_{22}\end{pmatrix}\in\Mat{2}{2}$. Determine si $\langle \cdot,\cdot\rangle $ es un producto interno.
%\end{ejer}

%\begin{ejer}\hspace{0pt} En el espacio vectorial $\R^2$, verifique si la función
%\[
 %   \func{\langle \cdot,\cdot\rangle }{\R^2\times\R^2}{\R}
%\]
%definida por
%\[
 %   \langle u,v\rangle  = u^{T}Av = (u_1,u_2)\begin{pmatrix}1 & 0\\ 2 & 0\end{pmatrix}\begin{pmatrix}v_1\\v_2\end{pmatrix}
%\]
%para $u^{T}=(u_1,u_2)$ y $v^{T}=(v_1,v_2)$, es un producto interno. 
%\end{ejer}

A continuación, se indican los productos internos usuales en los espacios vectoriales que hemos trabajado durante el curso.
\begin{enumerate}
\item En $(\R^n,+, \cdot, \R)$, para $x,y\in \R^n$:
\[
    \langle  x, y \rangle  = \dsum_{k=1}^{n}{x_k y_k}.
\]
\item En $(\C^n,+, \cdot, \C)$, para $x,y\in \C^n$:
\[
    \langle  x, y \rangle  = \dsum_{k=1}^{n}{x_k \overline{y_k}}.
\]
\item En $(\R_{n}[x], +, \cdot, \R)$, para $p(x),q(x)\in \R_{n}[x]$, si 
\[  
    p(x) = a_0 + a_1x + \cdots + a_nx^n
    \texty
    q(x) = b_0 + b_1x + \cdots + b_nx^n
\]
entonces
\[
    \langle  p(x), q(x) \rangle  = \dsum_{k=0}^{n}{a_k b_k}.
\]
\item En $(\Mat[\R]{n}{n},+, \cdot, \R)$, para $A,B\in \Mat[\R]{n}{n}$:
\[
    \langle  A, B \rangle  = \Tr(AB^{T}).
\]
%\item En $(\mathcal{C}([a,b]),+, \cdot, \R)$, para $f,g\in \mathcal{C}([a,b])$:
%\[
 %   \langle  f, g \rangle  = \int_{a}^{b}{f(x) g(x)\,dx}.
%\]
\end{enumerate}

\begin{teo}\label{te12:01}
    Sea $(E,+,\cdot,\K)$ un espacio vectorial provisto de producto interno, entonces:
    \begin{enumerate}
        \item 
            Para todo $u,v,w \in E$ se verifica que:
            \[\langle  u, v+w\rangle  = \langle  u,v\rangle  + \langle  u, w \rangle .\]
        \item 
            Para todo $u,v \in E$ y $\alpha \in \K$ se verifica que:
            \[\langle  u, \alpha v\rangle  =  \overline{\alpha}\langle  u,v\rangle .\]
        \item
            Para todo $u\in E$ se cumple que:
            \[\langle  u, 0 \rangle  =  \langle  0, u \rangle   = 0.\]
    \end{enumerate}
\end{teo}

\begin{advertencia}
    En caso de tratarse de un espacio vectorial sobre el campo de los números reales, la propiedad 2 del Teorema \ref{te12:01} implica
    \begin{enumerate}[start=2]
        \item
            Para todo $u,v \in E$ y $\alpha \in \R$
            \[\langle  u, \alpha v\rangle  =  \alpha\langle  u,v\rangle \]
    \end{enumerate}
\end{advertencia}

\begin{proof}[Demostración de la propiedad 2 del Teorema \ref{te12:01}.]\hspace{0pt} Sean $u,v\in E$ y $\alpha\in\K$, tenemos que
\begin{align*}
    \langle u,\alpha v\rangle  &= \overline{\langle \alpha v,u\rangle }\\
    &= \overline{\alpha\langle v,u\rangle }\\
    &= \overline{\alpha}\overline{\langle v,u\rangle }\\
    &= \overline{\alpha}\langle u,v\rangle .
\end{align*}
Lo que muestra al enunciado.
\end{proof}

\begin{ejer}\hspace{0pt} Sea $(E,+,\cdot,\R)$ un espacio vectorial provisto de producto interno, sean $u,v,w\in E$, demuestre que $\langle u - v,v - w\rangle  = \langle v,u - v\rangle  + \langle w,v-u\rangle .$
\end{ejer}

\begin{defi}[Vectores ortogonales]
Sean $(E,+,\cdot,\K)$ un espacio vectorial provisto de producto interno y $u, v \in E$. Se dice que $u$ y $v$ son ortogonales si $\langle  u, v \rangle  = 0$.
\end{defi}

\begin{ejem}\hspace{0pt}\label{Sem08:EjemOrt} En el espacio vectorial $\R^{3}$ con el producto interno usual, los vectores $u=(1,1,2)$ y $v=(2,0,-1)$ son ortogonales. En efecto, notemos que
\[
    \langle u,v\rangle  = \langle (1,1,2),(2,0,-1)\rangle  = (1)(2) + (1)(0) + (2)(-1) = 0.
\]
\end{ejem}

\begin{ejer}\hspace{0pt} En el espacio vectorial $\Mat{2}{2}$ con producto interno usual, considere a los vectores $A = \begin{pmatrix}-1&0\\-1&0\end{pmatrix}$ y $B = \begin{pmatrix}0&1\\1&1\end{pmatrix}$. Determine si $A$ y $B$ son ortogonales.
\end{ejer}

\begin{defi}[Norma]
Sea $(E,+,\cdot,\K)$ un espacio vectorial provisto de producto interno. Sea $u\in E$, la norma de $u$, denotada por $\|u \|$, está definida por:
\[
    \| u\| = \sqrt{\langle  u, u \rangle .}
\]
\end{defi}

\begin{ejem}\hspace{0pt} En el espacio vectorial $\mathbb{R}^{3}$ considere el vector $u=(2,3,-4)$ y halle su norma.\newline

Considerando el producto interno se tiene que:
\[
\| u\| = \sqrt{\langle  (2,3,-4), (2,3,-4) \rangle } = \sqrt{29}
\]
\end{ejem}

\begin{ejer}\hspace{0pt} En el espacio vectorial $\mathbb{R}^{2\times 2}$ considere la matriz $v=\begin{pmatrix}0&2\\1&-1\end{pmatrix}$ y halle su norma.
\end{ejer}

\begin{defi}[Distancia]
Sea $(E,+,\cdot,\K)$ un espacio vectorial con producto interno. La distancia en el espacio se define por
\[
    \funcion{d}{E\times E}{\R}{(u,v)}{\| u - v\|.}
\]
\end{defi}

\begin{ejem}\hspace{0pt} En el espacio vectorial $\R^{3}$ con el producto interno usual se tiene que la distancia en $\R^{3}$ es:
\[
    d(u, v) = \|u - v\| = \sqrt{(u_{1} - v_{1})^{2} + (u_{2} - v_{2})^{2} + (u_{3} - v_{3})^{2}}
\]
para todo $u = (u_{1}, u_{2}, u_{3}), v = (v_{1}, v_{2}, v_{3})\in\R^{3}$.
\end{ejem}


\begin{ejer}\hspace{0pt} En el espacio vectorial $\R^{3}$ considere los vectores:
\[u=(2,1,3)\]
\[v=(0,1,-2)\]
obtenga la distancia entre $u$ y $v$.

%En el espacio vectorial $\R_{2}[t]$ considere los vectores:
%\[p(t)=2+t+3t^2\]
%\[q(t)=t-2t^2\]
%obtenga la distancia entre $p$ y $q$.
\end{ejer}

%\begin{defi}[Norma]
%Sea $(E,+,\cdot,\K)$ un espacio vectorial provisto de producto interno. Sea %$u\in E$, la norma de $u$, denotada por $\|u \|$, está definida por:
%\[
%    \| u\| = \sqrt{\langle  u, u \rangle .}
%\]
%\end{defi}

%\begin{ejem}\hspace{0pt}\label{Sem08:EjemNorm} En el espacio vectorial %$\R^{3}$ con el producto interno usual, los vectores $u=(1,1,2)$ y %$v=(2,0,-1)$ tienen norma $\sqrt{6}$ y $\sqrt{5}$, respectivamente. En %efecto, notemos que
%\[
 %   \|u\| = \sqrt{\langle u,u\rangle } = \sqrt{\langle (1,1,2),(1,1,2)\rangle } = \sqrt{(1)^{2} + %(1)^{2} + (2)^{2}} = \sqrt{6}
%\]
%y
%\[
 %   \|v\| = \sqrt{\langle v,v\rangle } = \sqrt{\langle (2,0,-1),(2,0,-1)\rangle } = \sqrt{(2)^{2} %+ (0)^{2} + (-1)^{2}} = \sqrt{5}.
%\]
%\end{ejem}

%\begin{ejer}\hspace{0pt} En el espacio vectorial $\Mat{2}{2}$ con producto %interno, considere a los vectores $A = %\begin{pmatrix}-1&0\\-1&0\end{pmatrix}$ y $B = %\begin{pmatrix}0&1\\1&1\end{pmatrix}$. Determine sus normas.
%\end{ejer}

%\begin{defi}[Distancia]
%Sea $(E,+,\cdot,\K)$ un espacio vectorial con producto interno. La distancia %en el espacio se define por
%\[
 %   \funcion{d}{E\times E}{\R}{(u,v)}{\| u - v\|.}
%\]
%\end{defi}

%\begin{ejem}\hspace{0pt} En el espacio vectorial $\R^{3}$ con el producto %interno usual se tiene que la distancia en $\R^{3}$ es:
%\[
%    d(u, v) = \|u - v\| = \sqrt{(u_{1} - v_{1})^{2} + (u_{2} - v_{2})^{2} + %(u_{3} - v_{3})^{2}}
%\]
%para todo $u = (u_{1}, u_{2}, u_{3}), v = (v_{1}, v_{2}, v_{3})\in\R^{3}$. %Además, sean $u=(1,1,2)$ y $v=(2,0,-1)$ su distancia es $\sqrt{11}$. En %efecto, notemos que
%\[
 %   d(u, v) = \|u - v\| = \sqrt{(1 - 2)^{2} + (1 - 0)^{2} + (2 - (-1))^{2}} = %\sqrt{11}.
%\]
%\end{ejem}

%\begin{ejer}\hspace{0pt} En el espacio vectorial $\Mat{2}{2}$ con producto %interno, considere a los vectores $A = %\begin{pmatrix}-1&0\\-1&0\end{pmatrix}$ y $B = %\begin{pmatrix}0&1\\1&1\end{pmatrix}$. Determine su distancia.
%\end{ejer}

\begin{teo}[Teorema de Pitágoras]
    Sea $(E,+,\cdot,\K)$ un espacio vectorial con producto interno.
    Si $u,v$ son vectores ortogonales de $E$, entonces
    \[
        \| u + v \|^2 = \|u\|^2 + \|v\|^2.
    \]
\end{teo}
\begin{proof}
    Sea $(E,+,\cdot,\K)$ un espacio vectorial con producto interno.
    Tal que  $u,v$ son vectores ortogonales de $E$, por definición de norma sabemos que 
    \[
    \| u + v \|^2 =\langle u+v,u+v\rangle.
    \]
    Por teorema \ref{te12:01} se tiene que 
    \begin{align*}
         \| u + v \|^2 &=\langle u+v,u+v\rangle\\
         &=\langle u+v,u\rangle+\langle u+v,v\rangle\\
         &=\langle u,u\rangle+\langle v,u\rangle+\langle u,v\rangle+\langle v,v\rangle.\\
    \end{align*}
    Además como $u,v$ son vectores ortogonales sabemos que $\langle v,u\rangle = \langle u,v\rangle =0 $, así reemplazando en la anterior ecuación tenemos que 
    \[
    \| u + v \|^2=\langle u,u\rangle+\langle v,v\rangle.
    \]
    Finalmente, usando la definición de norma se sigue que
    \[
    \| u + v \|^2 = \|u\|^2 + \|v\|^2.
    \]
\end{proof}

\begin{ejem}\hspace{0pt}\label{Sem08:EjemPit} En el espacio vectorial $\R^{3}$ con el producto interno usual, los vectores $u=(1,1,2)$ y $v=(2,0,-1)$ son ortogonales y son tales que 
\[
    \| u + v \|^2 = \|u\|^2 + \|v\|^2.
\]
En efecto, por el ejemplo de vectores ortogonales, $u$ y $v$ son ortogonales y por el ejemplo de norma, sus normas son tales que $\|u\|^{2}=6$ y $\|v\|^{2}=5$, además,
\[
    \|u + v\|^{2} = \|(3,1,1)\|^{2} = (3)^{2} + (1)^{2} + (1)^{2} = 11=\|u\|^{2}+\|v\|^{2}.
\]
Lo que verifica el resultado.
\end{ejem}

\begin{ejem}
    Sea $(E,+,\cdot,\K)$ un espacio vectorial con producto interno.
    Si $u,v$ son vectores ortogonales de $E$, tales que 
    \[
    \|u\|=\|v\|=1
    \]
    Hallar $\|u+v\|$.

    Por hipótesis sabemos que $u,v$ son ortogonales, por tanto 
    \begin{align*}
        \|u+v\|^2&=\|u\|^2+\|v\|^2\\
        &=1^2+1^2\\
        &=1+1\\
        &=2
    \end{align*}
    Finalmente, 
    \[
    \|u+v\|=\sqrt{2}
    \]
\end{ejem}
\begin{ejer}
    Sea $(E,+,\cdot,\K)$ un espacio vectorial con producto interno.
    Si $u,v$ son vectores ortogonales de $E$, tales que 
    \[
    \|u\|=\|v\|=2
    \]
    Hallar $\|u-v\|$.
\end{ejer}
%\begin{teo}[Desigualdad de Cauchy-Schwartz]
 %   Sea $(E,+,\cdot,\K)$ un espacio vectorial provisto de producto interno.
  %  Para todo $u,v\in E$ se cumple que
   % \[
    %    | \langle  u , v \rangle  | \leq \|u\|  \|v\|.
    %\]
%\end{teo}
%\begin{advertencia}
 %   Esta desigualdad es una igualdad si y solo si $u$ es un múltiplo escalar de $v$.
%\end{advertencia}
%\comentario{Opción 1) para demostrar en CD.}
\begin{comment}
\begin{proof}%Demostración (opcion 1)
Sean $u,v \in E$. 

Si $v=0$ entonces la desigualdad se cumple pues ambos lados de la desigualdad son cero.

Si $ v\neq 0$, considere la descomposición ortogonal
\begin{equation*}\label{eq12:03}
    u = \dfrac{\langle  u , v \rangle }{\|v\|^2} v + w
\end{equation*}
donde $w\in E$ es ortogonal a $v$. Aplicando el Teorema de Pitágoras a la desigualdad a \eqref{eq12:03} se obtiene
\begin{align*}
    \|u\|^2 
&= \left| \dfrac{\langle  u,v \rangle }{\|v\|^2} v \right\| + \|w\|^2 \\
&= \dfrac{\|\langle  u , v \rangle  |^2 \|v\|^2}{\|v\|^4 + \|w\|^2}\\
&= \dfrac{|\langle  u,v \rangle  | }{\|v\|^2} + \|w\|^2 \\
&\geq \dfrac{| \langle  u,v \rangle  |^2}{ \|v\|^2 }
\end{align*}
luego tomando la raíz cuadrada a cada lado se obtiene la desigualdad de Cauchy-Schwartz. 
\end{proof}
\end{comment}

\end{document}
